\documentclass[border=10pt,tikz]{standalone}
\usepackage{tikz}
\usetikzlibrary{positioning,arrows.meta,shapes,fit,backgrounds,calc}
\usepackage{kotex}

% JTAG TAP 상태 머신 (16 상태).
% 상태 이름: OpenOCD src/jtag/interface.c tap_name_mapping[]
% 전이: OpenOCD src/jtag/interface.c tap_state_transition()
% 안정 상태 6개: OpenOCD doc/openocd.texi "TAP state names"

\begin{document}
% _design.tex — shared TikZ design system for all blog diagrams
% Inputs nothing on its own — meant to be \input{} from each diagram .tex
% AFTER \begin{document}, BEFORE \begin{tikzpicture}.
%
% Provides a unified visual language across GoF / DSA / EMC++ / Beautiful series:
%   - Colors: muted, accessible, color-coded by role
%   - Node styles: consistent sizing, rounded, subtle borders
%   - Edge styles: UML-style inheritance, aggregation, plus dashed alternates
%
% Usage:
%   \documentclass[border=10pt,tikz]{standalone}
%   \usepackage{tikz}
%   \usetikzlibrary{positioning,arrows.meta,shapes,fit,backgrounds}
%   \begin{document}
%   \input{../_design.tex}     % adjust relative path
%   \begin{tikzpicture}[blog]  % apply 'blog' style
%   ...

\definecolor{absfill}{RGB}{220,232,247}   % cool blue
\definecolor{absbord}{RGB}{72,118,182}
\definecolor{confill}{RGB}{249,222,222}   % warm red
\definecolor{conbord}{RGB}{197,93,93}
\definecolor{impfill}{RGB}{223,238,222}   % soft green
\definecolor{impbord}{RGB}{99,154,93}
\definecolor{clifill}{RGB}{251,239,205}   % pale yellow
\definecolor{clibord}{RGB}{200,167,72}
\definecolor{hifill} {RGB}{253,224,200}   % accent orange
\definecolor{hibord} {RGB}{210,135,68}
\definecolor{nodefill}{RGB}{250,250,250}
\definecolor{nodebord}{RGB}{120,120,120}
\definecolor{edgec}  {RGB}{80,80,80}
\definecolor{edgesc} {RGB}{145,145,145}
\definecolor{groupbord}{RGB}{200,200,200}

\tikzset{
  blog/.style={
    font=\sffamily\small,
    every node/.style={font=\sffamily\small}
  },
  % --- node roles ---
  cls/.style={
    rectangle, rounded corners=2pt,
    draw=nodebord, line width=0.4pt,
    minimum height=8mm, minimum width=28mm,
    inner sep=3pt, fill=nodefill, align=center
  },
  abs/.style={cls, draw=absbord, fill=absfill, font=\sffamily\itshape\small},
  con/.style={cls, draw=conbord, fill=confill},
  imp/.style={cls, draw=impbord, fill=impfill},
  cli/.style={cls, draw=clibord, fill=clifill},
  hi/.style={cls, draw=hibord, fill=hifill, font=\sffamily\bfseries\small},
  % --- circle (for trees / graph nodes) ---
  vx/.style={
    circle, draw=absbord, line width=0.4pt,
    minimum size=8mm, fill=absfill, inner sep=0pt
  },
  vxc/.style={vx, draw=conbord, fill=confill},
  vxh/.style={vx, draw=hibord, fill=hifill, font=\sffamily\bfseries},
  % --- decision diamond ---
  q/.style={
    diamond, draw=clibord, line width=0.4pt, fill=clifill,
    inner sep=2pt, aspect=2.4, align=center
  },
  % --- edges ---
  inh/.style={                                  % UML inheritance
    -{Triangle[length=3mm,width=3mm,open]},
    draw=edgec, line width=0.5pt
  },
  agg/.style={                                  % UML aggregation
    {Diamond[length=2.5mm,width=2.5mm,open]}-{Stealth[length=2.2mm]},
    draw=edgec, line width=0.45pt
  },
  uses/.style={                                 % directed reference
    -{Stealth[length=2.2mm]}, draw=edgec, line width=0.5pt
  },
  arr/.style={                                  % plain arrow (graphs/trees)
    -{Stealth[length=2mm]}, draw=edgec, line width=0.45pt
  },
  edge/.style={                                 % undirected (trees)
    draw=edgec, line width=0.45pt
  },
  alt/.style={                                  % dashed: similar/alt
    -{Stealth[length=2mm]}, dashed, draw=edgesc, line width=0.45pt
  },
  ret/.style={                                  % dashed creates/returns
    -{Stealth[length=2mm]}, dashed, draw=edgesc, line width=0.45pt
  },
  thr/.style={                                  % threading link (red dashed)
    ->, dashed, draw=conbord, line width=0.5pt
  },
  % --- group / region ---
  group/.style={
    rectangle, rounded corners=4pt,
    draw=groupbord, line width=0.4pt,
    inner sep=10pt, fill=black!2
  },
  glabel/.style={font=\sffamily\bfseries\small, anchor=north west, inner sep=2pt}
}

% Korean font convenience macro — included by every Hangul diagram
\providecommand{\KOR}{}

% _design-state.tex — TikZ state machine / FSM styles
%
% UML state diagram, protocol state machine, hardware FSM, kernel state graph
% 등에 공통. Mermaid의 state diagram은 *수동 정렬*이 보통 더 깔끔하므로
% TikZ로 직접 그리는 것이 표준 (CLAUDE.md 규칙).
%
% Companion to _design.tex — same palette and ink colors.
%
% Usage:
%   \documentclass[border=10pt,tikz]{standalone}
%   \usepackage{tikz}
%   \usetikzlibrary{positioning,arrows.meta,shapes,fit,calc,decorations.pathreplacing}
%   \begin{document}
%   \input{../_design.tex}        % core palette
%   \input{../_design-state.tex}  % state-specific
%   \begin{tikzpicture}[blog, state]
%     \node[initial]            (i)   at (-2, 0)  {};
%     \node[stateBox]           (s1)  at ( 0, 0)  {Idle};
%     \node[stateBox]           (s2)  at ( 3, 0)  {Running};
%     \node[stateBox, accent]   (s3)  at ( 6, 0)  {Done};
%     \node[final]              (f)   at ( 8, 0)  {};
%
%     \draw[trans] (i)  -- (s1);
%     \draw[trans] (s1) -- node[transLabel]{start} (s2);
%     \draw[trans] (s2) -- node[transLabel]{finish} (s3);
%     \draw[trans] (s3) -- (f);
%
%     % self transition
%     \draw[trans] (s2) edge[loop above] node[transLabel]{tick} (s2);
%
%     % guard / action label
%     \draw[trans] (s1) to[bend right=30] node[transLabel]{stop\,/\,cleanup} (s3);
%
%     % choice (junction)
%     \node[choice] (c) at (4.5, -2) {};
%   \end{tikzpicture}

\tikzset{
  state/.style={
    every node/.append style={font=\sffamily\small},
  },

  % ─── State (rounded rectangle, default) ─────────────────
  stateBox/.style={
    rectangle, rounded corners=4pt,
    draw=nodebord, line width=0.5pt,
    minimum height=8mm, minimum width=22mm,
    inner sep=4pt, fill=nodefill, align=center,
  },
  stateBox/.append style={font=\sffamily\small},
  % Role-based variants — match _design.tex palette
  stateAbs/.style ={stateBox, draw=absbord, fill=absfill, font=\sffamily\itshape\small},
  stateCon/.style ={stateBox, draw=conbord, fill=confill},
  stateImp/.style ={stateBox, draw=impbord, fill=impfill},
  stateOk/.style  ={stateBox, draw=okbord,  fill=okfill},
  stateWarn/.style={stateBox, draw=warnbord, fill=warnfill, font=\sffamily\bfseries\small},
  statePrp/.style ={stateBox, draw=prpbord, fill=prpfill},   % protocol / interface
  stateNeut/.style={stateBox, draw=neutbord, fill=neutfill},
  accent/.style   ={draw=hibord, fill=hifill},               % apply on top (e.g., \node[stateBox, accent])

  % Compact / large variants
  stateSm/.style ={stateBox, minimum height=6mm, minimum width=16mm, font=\sffamily\footnotesize, inner sep=2pt},
  stateLg/.style ={stateBox, minimum height=11mm, minimum width=30mm},

  % ─── Initial / final / terminal ─────────────────────────
  initial/.style={
    circle, fill=inkc, draw=inkc, line width=0.5pt,
    minimum size=3.5mm, inner sep=0pt,
  },
  final/.style={                                              % bullseye (outer ring + inner dot)
    circle, draw=inkc, line width=0.5pt,
    minimum size=5mm, inner sep=0pt,
    path picture={
      \fill[inkc] (path picture bounding box.center) circle (1.2mm);
    },
  },
  terminate/.style={                                          % X (terminate pseudo-state)
    cross out, draw=warnbord, line width=0.8pt,
    minimum size=3mm, inner sep=0pt,
  },

  % ─── Choice / junction / fork-join ──────────────────────
  choice/.style={                                             % diamond decision
    diamond, draw=clibord, line width=0.45pt, fill=clifill,
    aspect=1, inner sep=1pt, minimum size=6mm,
    font=\sffamily\scriptsize, align=center,
  },
  junction/.style={                                           % small filled black circle (UML junction)
    circle, fill=inkc, draw=inkc, line width=0.4pt,
    minimum size=2.5mm, inner sep=0pt,
  },
  forkbar/.style={                                            % horizontal bar (fork / join)
    rectangle, fill=inkc, draw=inkc, line width=0.4pt,
    minimum width=20mm, minimum height=1.2mm, inner sep=0pt,
  },
  forkbarV/.style={forkbar, minimum width=1.2mm, minimum height=20mm},

  % ─── Transitions ────────────────────────────────────────
  trans/.style={
    -{Stealth[length=2.4mm]}, draw=edgec, line width=0.5pt,
  },
  transAlt/.style={                                           % dashed (lower priority / alt)
    -{Stealth[length=2.4mm]}, dashed, draw=edgesc, line width=0.5pt,
  },
  transFocus/.style={                                         % highlighted hot path
    -{Stealth[length=2.6mm]}, draw=focusc, line width=0.7pt,
  },
  transError/.style={                                         % error transition
    -{Stealth[length=2.4mm]}, draw=warnbord, line width=0.55pt,
  },
  % Label on transition — event[guard]/action 형식
  transLabel/.style={
    font=\sffamily\footnotesize, text=inkc,
    fill=white, fill opacity=0.85, text opacity=1,
    inner sep=2pt,
  },
  transLabelMute/.style={transLabel, text=mutec},

  % ─── Internal action listing (inside state box) ─────────
  internalAction/.style={
    font=\sffamily\footnotesize, text=inkc, anchor=north west, inner sep=2pt,
  },
  separator/.style={                                          % thin separator inside composite state
    draw=mutec, line width=0.3pt,
  },

  % ─── Composite / region (containing sub-states) ─────────
  composite/.style={
    rectangle, rounded corners=6pt,
    draw=nodebord, line width=0.5pt,
    inner sep=10pt, fill=neutfill!50,
  },
  compositeTitle/.style={
    font=\sffamily\bfseries\small, text=inkc,
    fill=neutfill, draw=neutbord, line width=0.3pt,
    inner sep=2pt, rounded corners=2pt,
    anchor=north west,
  },
  region/.style={                                             % parallel region inside composite
    rectangle, dashed,
    draw=mutec, line width=0.35pt,
    inner sep=8pt,
  },
}


\begin{tikzpicture}[blog, state,
    st/.style={stateBox, minimum width=24mm, minimum height=8mm, font=\sffamily\small},
    stable/.style={st, draw=absbord, fill=absfill, font=\sffamily\small\bfseries},
    tms/.style={font=\sffamily\footnotesize, text=inkc, inner sep=1.5pt},
    ln/.style={draw=edgec, line width=0.5pt}
]

% ---- columns: control (x=0), DR (x=5), IR (x=11.2)
\node[stable] (reset) at (0, 0)    {RESET};
\node[stable] (idle)  at (0, -2)   {RUN/IDLE};

\node[st]     (drs)  at (5, -2)    {DRSELECT};
\node[st]     (drc)  at (5, -3.5)  {DRCAPTURE};
\node[stable] (drsh) at (5, -5)    {DRSHIFT};
\node[st]     (dre1) at (5, -6.5)  {DREXIT1};
\node[stable] (drp)  at (5, -8)    {DRPAUSE};
\node[st]     (dre2) at (5, -9.5)  {DREXIT2};
\node[st]     (dru)  at (5, -11)   {DRUPDATE};

\node[st]     (irs)  at (11.2, -2)    {IRSELECT};
\node[st]     (irc)  at (11.2, -3.5)  {IRCAPTURE};
\node[stable] (irsh) at (11.2, -5)    {IRSHIFT};
\node[st]     (ire1) at (11.2, -6.5)  {IREXIT1};
\node[stable] (irp)  at (11.2, -8)    {IRPAUSE};
\node[st]     (ire2) at (11.2, -9.5)  {IREXIT2};
\node[st]     (iru)  at (11.2, -11)   {IRUPDATE};

% ---- control column
% RESET self-loop (TMS=1)
\draw[trans] ([yshift=1.5mm]reset.west) -- ++(-0.6,0) |- ([yshift=-1.5mm]reset.west);
\node[tms, anchor=east] at ([xshift=-6.5mm]reset.west) {1};
\draw[trans] (reset) -- node[tms, right]{0} (idle);
% RUN/IDLE self-loop (TMS=0)
\draw[trans] ([yshift=1.5mm]idle.west) -- ++(-0.6,0) |- ([yshift=-1.5mm]idle.west);
\node[tms, anchor=east] at ([xshift=-6.5mm]idle.west) {0};
\draw[trans] (idle) -- node[tms, above]{1} (drs);

% ---- select row
\draw[trans] (drs) -- node[tms, above]{1} (irs);
\draw[trans] (irs.north) |- node[tms, pos=0.25, right]{1} (reset.east);

% ---- one column of the scan path (DR or IR); #1 = prefix, #2 = x
\newcommand{\scancol}[2]{%
  \draw[trans] (#1s) -- node[tms, right]{0} (#1c);
  \draw[trans] (#1c) -- node[tms, right]{0} (#1sh);
  \draw[trans] (#1sh) -- node[tms, right]{1} (#1e1);
  \draw[trans] (#1e1) -- node[tms, right]{0} (#1p);
  \draw[trans] (#1p) -- node[tms, right]{1} (#1e2);
  \draw[trans] (#1e2) -- node[tms, right]{1} (#1u);
  % CAPTURE -> EXIT1 (TMS=1), left lane 1
  \draw[trans] (#1c.west) -- (#2-1.8, -3.5) -- node[tms, left, pos=0.5]{1} (#2-1.8, -6.35) -- ([yshift=1.5mm]#1e1.west);
  % EXIT1 -> UPDATE (TMS=1), left lane 2
  \draw[trans] ([yshift=-1.5mm]#1e1.west) -- (#2-2.4, -6.65) -- node[tms, left, pos=0.5]{1} (#2-2.4, -11) -- (#1u.west);
  % EXIT2 -> SHIFT (TMS=0), right lane
  \draw[trans] (#1e2.east) -- (#2+2.3, -9.5) -- node[tms, right, pos=0.5]{0} (#2+2.3, -5.25) -- ([yshift=-2.5mm]#1sh.east);
  % SHIFT / PAUSE self-loops (TMS=0)
  \draw[trans] ([yshift=3mm]#1sh.east) -- ++(0.5,0) |- ([yshift=0.5mm]#1sh.east);
  \node[tms, anchor=west] at ([xshift=5.5mm, yshift=1.75mm]#1sh.east) {0};
  \draw[trans] ([yshift=3mm]#1p.east) -- ++(0.5,0) |- ([yshift=0.5mm]#1p.east);
  \node[tms, anchor=west] at ([xshift=5.5mm, yshift=1.75mm]#1p.east) {0};
}
\scancol{dr}{5}
\scancol{ir}{11.2}

% ---- UPDATE -> DRSELECT (TMS=1): both UPDATE states share one return lane
\draw[ln] ([xshift=-4mm]iru.south) -- ++(0,-0.4) -| (7.9, -11);
\draw[trans] (dru.east) -- (7.9, -11) -- node[tms, right, pos=0.5]{1} (7.9, -2.75) -- (5.8, -2.75) -- (5.8, -2.4);
\node[tms, anchor=west] at ([xshift=-3.6mm, yshift=-2mm]iru.south) {1};

% ---- UPDATE -> RUN/IDLE (TMS=0): shared bottom lane
\draw[ln] (dru.south) -- (5, -12.3);
\node[tms, anchor=west] at ([xshift=0.5mm, yshift=-2mm]dru.south) {0};
\draw[trans] ([xshift=4mm]iru.south) -- (11.6, -12.3) -- (0, -12.3) -- (idle.south);
\node[tms, anchor=west] at ([xshift=4.5mm, yshift=-2mm]iru.south) {0};

% ---- legend
\node[stable, minimum width=8mm, minimum height=5mm] (lg1) at (0.4, -13.4) {};
\node[anchor=west, font=\sffamily\footnotesize, text=inkc] at ([xshift=2mm]lg1.east)
  {TMS를 고정해도 머무는 안정 상태 6개};
\node[st, minimum width=8mm, minimum height=5mm] (lg2) at (7.2, -13.4) {};
\node[anchor=west, font=\sffamily\footnotesize, text=inkc] at ([xshift=2mm]lg2.east)
  {다음 TCK에서 넘어가는 상태};
\node[anchor=west, font=\sffamily\footnotesize, text=mutec] at (-0.4, -14.2)
  {화살표 숫자는 TMS 값입니다. 상태 이름은 OpenOCD 표기(RUN/IDLE = Run-Test/Idle)입니다.};

\end{tikzpicture}
\end{document}
